% Copyright (c) 2024
% Public domain.

\chapter{主要符号对照表}

本附录列出论文中使用的英文缩写及其对应的全称和中文含义，按字母顺序排列。

\begin{longtable}{llp{6cm}}
\toprule
\textbf{缩写} & \textbf{英文全称} & \textbf{中文含义} \\
\midrule
\endhead
\bottomrule
\endfoot
API      & Application Programming Interface           & 应用程序接口 \\
ASR      & Attack Success Rate                         & 攻击成功率 \\
CAI      & Constitutional AI                           & 宪法人工智能 \\
CLM      & Causal Language Modeling                    & 因果语言建模 \\
CoT      & Chain-of-Thought                            & 思维链 \\
DPO      & Direct Preference Optimization              & 直接偏好优化 \\
GPU      & Graphics Processing Unit                    & 图形处理器 \\
IPO      & Identity Preference Optimization            & 身份偏好优化 \\
KTO      & Kahneman-Tversky Optimization               & Kahneman-Tversky优化 \\
LLM      & Large Language Model                        & 大语言模型 \\
MLM      & Masked Language Modeling                    & 掩码语言建模 \\
PEFT     & Parameter-Efficient Fine-Tuning             & 参数高效微调 \\
PPO      & Proximal Policy Optimization                & 近端策略优化 \\
RAO      & Relative Advantage Optimization             & 相对优势优化 \\
RAG      & Retrieval-Augmented Generation              & 检索增强生成 \\
RLHF     & Reinforcement Learning from Human Feedback  & 人类反馈强化学习 \\
RLVR     & Rule Verification Reinforcement Learning    & 规则验证强化学习 \\
SafetyCoT& Safety Chain-of-Thought                     & 安全思维链 \\
SFT      & Supervised Fine-Tuning                      & 监督微调 \\
VCoT     & Visual Chain-of-Thought                     & 视觉思维链 \\
VLM      & Vision-Language Model                       & 视觉语言模型 \\
\end{longtable}

% vim:ts=4:sw=4
